\documentclass[11pt,letterpaper]{article}

% ==========================================================
%  RiskGate - Actualizacion del modelo de riesgo 2026
% ==========================================================

\usepackage[spanish,es-noshorthands]{babel}
\usepackage[utf8]{inputenc}
\usepackage[T1]{fontenc}
\usepackage{lmodern}
\usepackage[margin=1in]{geometry}
\usepackage{microtype}
\usepackage{amsmath}
\usepackage{amssymb}
\usepackage{array}
\usepackage{booktabs}
\usepackage{longtable}
\usepackage{enumitem}
\usepackage{xcolor}
\usepackage{listings}
\usepackage{hyperref}

\hypersetup{
  colorlinks=true,
  linkcolor=blue,
  urlcolor=blue,
  citecolor=blue,
  pdftitle={RiskGate - Actualizacion del modelo de riesgo 2026},
  pdfauthor={RiskGate}
}

\definecolor{rgcode}{gray}{0.96}
\lstdefinestyle{riskgate}{
  basicstyle=\ttfamily\small,
  breaklines=true,
  frame=single,
  columns=fullflexible,
  backgroundcolor=\color{rgcode},
  keepspaces=true,
  showstringspaces=false
}
\lstset{style=riskgate}

\newcommand{\RiskGate}{\textsc{RiskGate}}
\newcommand{\MXN}{\ensuremath{\mathrm{MXN}}}
\newcommand{\clp}{\operatorname{clamp}}
\setlist[itemize]{leftmargin=1.3em}
\setlist[enumerate]{leftmargin=1.5em}

\title{\textbf{\RiskGate}\\[2mm]
\large Actualizaci\'on del modelo de riesgo, probabilidad y sentimiento}
\author{Proyecto \RiskGate{} --- aplicaci\'on local-first}
\date{Mayo de 2026}

\begin{document}

\maketitle

\begin{abstract}
\noindent
Este documento resume las modificaciones recientes incorporadas a \RiskGate{}:
r\'egimen de mercado compuesto, jerarqu\'ia de probabilidad de ganancia basada en
opciones, dimensionamiento Kelly sin usar el puntaje como probabilidad, capa de
riesgo por eventos, validaciones espec\'ificas por estrategia y reemplazo del
sentimiento de Reddit por sentimiento de noticias de Yahoo Finance. La
aplicaci\'on permanece local-first, de solo lectura frente a fuentes externas,
sin credenciales de broker, sin colocaci\'on de \'ordenes y sin automatizaci\'on de
trading.
\end{abstract}

\tableofcontents
\newpage

% ==========================================================
\section{Resumen ejecutivo}
% ==========================================================

La actualizaci\'on corrige cinco debilidades del modelo anterior:

\begin{enumerate}[noitemsep]
  \item El multiplicador de r\'egimen ya no depende solo de SPY.
  \item La probabilidad de ganancia (PoP) ya no se infiere como
        \texttt{score / 100}.
  \item El tama\~no por Kelly usa una PoP derivada de opciones o se desactiva.
  \item Los eventos de earnings y calendario macro pueden advertir, reducir
        tama\~no, limitar la decisi\'on o rechazar estructuras severas.
  \item El sentimiento opcional ahora se basa en noticias de Yahoo Finance,
        no en Reddit.
\end{enumerate}

El flujo final del motor de riesgo queda organizado as\'i:

\begin{lstlisting}
1. Validar campos obligatorios.
2. Construir regimen de mercado compuesto.
3. Calcular puntajes base: mercado, tecnico, opciones y plan.
4. Aplicar pesos adaptativos.
5. Ejecutar analitica de opciones y Monte Carlo cuando haya datos.
6. Estimar PoP con jerarquia de fuentes.
7. Calcular Kelly desde PoP.
8. Evaluar riesgo por eventos.
9. Evaluar checks especificos de estrategia.
10. Aplicar overlay de Yahoo Finance Sentiment si esta activo.
11. Convertir score a decision.
12. Aplicar hard stops y caps.
13. Calcular riesgo sugerido con todos los multiplicadores.
\end{lstlisting}

% ==========================================================
\section{Nuevos m\'odulos}
% ==========================================================

\begin{longtable}{p{0.33\linewidth} p{0.57\linewidth}}
\toprule
\textbf{Archivo} & \textbf{Responsabilidad} \\
\midrule
\texttt{lib/risk/market-regime.ts} & Calcula el r\'egimen compuesto con SPY,
QQQ, amplitud de mercado y VIX. \\
\texttt{lib/options/pop.ts} & Estima PoP con jerarqu\'ia de fuentes y calcula
Kelly fraccional desde PoP. \\
\texttt{lib/risk/event-risk.ts} & Eval\'ua cruces de earnings y eventos macro,
produce advertencias, caps, hard stops y multiplicador de riesgo. \\
\texttt{lib/risk/strategy-checks.ts} & Valida reglas propias de long call, long
put, debit spreads, protective put y collar. \\
\texttt{lib/sentiment/yahoo/} & Recolecta, limpia, clasifica y aplica overlay
de sentimiento de noticias de Yahoo Finance. \\
\texttt{data/macro-events.json} & Calendario macro local editable. \\
\texttt{scripts/migrate-yahoo-sentiment-and-risk-overlays.mjs} & Migraci\'on
SQLite local para settings nuevos y cach\'e de documentos Yahoo. \\
\bottomrule
\end{longtable}

% ==========================================================
\section{R\'egimen de mercado compuesto}
% ==========================================================

Antes, el multiplicador de r\'egimen usaba principalmente la tendencia de SPY.
Ahora se calcula un \'indice compuesto normalizado:

\[
C \in [-1,1].
\]

El r\'egimen combina cuatro componentes:

\begin{longtable}{p{0.22\linewidth} p{0.18\linewidth} p{0.48\linewidth}}
\toprule
\textbf{Componente} & \textbf{Peso} & \textbf{Idea} \\
\midrule
SPY & 30\% & Tendencia del mercado amplio. \\
QQQ & 25\% & Tendencia de crecimiento/Nasdaq. \\
Breadth & 25\% & Proporci\'on de universo/watchlist arriba de SMA50. \\
VIX & 20\% & Volatilidad impl\'icita de mercado, invertida. \\
\bottomrule
\end{longtable}

Si falta un componente, el modelo normaliza los pesos disponibles:

\[
C = \frac{\sum_i s_i w_i}{\sum_i w_i\ \text{disponible}}.
\]

La confianza es el peso total disponible:

\[
\text{confianza} = \sum_i w_i\ \text{disponible}.
\]

\subsection{Etiquetas}

\begin{longtable}{p{0.35\linewidth} p{0.45\linewidth}}
\toprule
\textbf{Rango} & \textbf{Etiqueta} \\
\midrule
$C \ge 0.60$ & \texttt{strong\_bullish} \\
$0.25 \le C < 0.60$ & \texttt{bullish} \\
$-0.25 < C < 0.25$ & \texttt{neutral} \\
$-0.60 < C \le -0.25$ & \texttt{bearish} \\
$-0.80 < C \le -0.60$ & \texttt{strong\_bearish} \\
$C \le -0.80$ & \texttt{stress} \\
\bottomrule
\end{longtable}

Un VIX mayor o igual a 35, o un salto diario igual o mayor a 15\%, puede forzar
\texttt{stress}.

\subsection{Multiplicador de r\'egimen}

El multiplicador $G$ ahora depende de la etiqueta compuesta y de la direcci\'on
del trade.

\begin{longtable}{p{0.24\linewidth} p{0.20\linewidth} p{0.20\linewidth} p{0.20\linewidth}}
\toprule
\textbf{R\'egimen} & \textbf{Bullish} & \textbf{Bearish} & \textbf{Neutral/Hedge} \\
\midrule
\texttt{strong\_bullish} & 1.00 & 0.25 & 0.75 \\
\texttt{bullish} & 0.85 & 0.35 & 0.85 \\
\texttt{neutral} & 0.50 & 0.50 & 1.00 \\
\texttt{bearish} & 0.35 & 0.85 & 0.85 \\
\texttt{strong\_bearish} & 0.25 & 1.00 & 0.75 \\
\texttt{stress} & 0.15 & 0.75 & 0.50 \\
\bottomrule
\end{longtable}

Si la confianza es menor a 50\%, $G$ se limita a 0.50 y se muestra una
advertencia de baja confianza.

% ==========================================================
\section{PoP y Kelly}
% ==========================================================

El cambio m\'as importante de probabilidad es que el puntaje del modelo ya no es
tratado como probabilidad. Se elimin\'o la l\'ogica activa equivalente a:

\begin{lstlisting}
p = totalScore / 100
\end{lstlisting}

La nueva jerarqu\'ia de fuentes es:

\begin{enumerate}[noitemsep]
  \item \textbf{Monte Carlo}: prioridad para estructuras multi-leg con simulaci\'on
        v\'alida.
  \item \textbf{Lognormal con IV real}: usa spot, strike, DTE e IV.
  \item \textbf{Proxy por delta}: usa delta absoluta o delta neta aproximada.
  \item \textbf{Proxy por IV}: solo si falta IV real pero existe proxy calibrado.
  \item \textbf{No disponible}: PoP nula y Kelly desactivado.
\end{enumerate}

El resultado incluye:

\begin{lstlisting}
{
  pop: number | null,
  source: "monte_carlo" | "lognormal_iv" | "delta_proxy" |
          "iv_proxy" | "unavailable",
  confidence: "high" | "medium" | "low",
  warnings: string[],
  inputsUsed: { iv, delta, dte, spot, strike, monteCarloPaths }
}
\end{lstlisting}

\subsection{Kelly desde PoP}

Con $p$ como PoP y $b$ como recompensa/riesgo:

\[
f^* = \max\left(0,\frac{p b - (1-p)}{b}\right).
\]

\RiskGate{} usa Kelly fraccional al 25\%:

\[
f_{\text{RG}} = 0.25 f^*.
\]

El tope de riesgo por Kelly es:

\[
R_{\text{Kelly}} = V_{\text{portafolio}} f_{\text{RG}}.
\]

Si PoP no est\'a disponible:

\begin{itemize}[noitemsep]
  \item Kelly se desactiva.
  \item El riesgo sugerido se vuelve cero.
  \item La decisi\'on queda limitada como m\'aximo a \texttt{watchlist}.
  \item Se muestra la advertencia correspondiente.
\end{itemize}

Si PoP tiene confianza baja, la decisi\'on se limita como m\'aximo a
\texttt{approved\_small}.

% ==========================================================
\section{Riesgo por eventos}
% ==========================================================

El nuevo overlay de eventos no reemplaza el puntaje base. Act\'ua despu\'es del
puntaje y antes del dimensionamiento final. Puede:

\begin{itemize}[noitemsep]
  \item agregar advertencias;
  \item reducir el multiplicador de riesgo;
  \item limitar la decisi\'on a \texttt{watchlist} o \texttt{approved\_small};
  \item forzar \texttt{reject} ante hard stops.
\end{itemize}

El archivo \texttt{data/macro-events.json} contiene eventos macro locales como
CPI, FOMC, PCE y NFP. El calendario es editable y no se hardcodean fechas futuras
en TypeScript.

\subsection{Reglas principales}

\begin{itemize}[noitemsep]
  \item Si el vencimiento cruza earnings, se advierte sobre gap risk e IV crush.
  \item Un debit spread que cruza earnings puede quedar limitado a
        \texttt{approved\_small}.
  \item Exposici\'on naked short cruzando earnings genera hard stop.
  \item Eventos macro cercanos reducen el multiplicador a 0.75 o 0.50.
  \item Evento macro en r\'egimen \texttt{stress} limita la decisi\'on a
        \texttt{watchlist}.
\end{itemize}

El resultado del m\'odulo tiene esta forma conceptual:

\begin{lstlisting}
{
  eventRiskLevel: "none" | "low" | "medium" | "high",
  crossesEarnings: boolean,
  crossesMacroEvent: boolean,
  events: RiskEvent[],
  eventWarnings: string[],
  hardStops: string[],
  decisionCap?: "watchlist" | "approved_small",
  riskMultiplier: number
}
\end{lstlisting}

% ==========================================================
\section{Checks espec\'ificos por estrategia}
% ==========================================================

El modelo ahora valida la estructura de cada estrategia, no solo reglas
generales de calidad.

\begin{longtable}{p{0.30\linewidth} p{0.58\linewidth}}
\toprule
\textbf{Estrategia} & \textbf{Checks principales} \\
\midrule
\texttt{long\_call} & DTE corto, IV rank alto, theta elevada, reward/risk
d\'ebil y alineaci\'on con r\'egimen compuesto. \\
\texttt{long\_put} & Reglas an\'alogas a long call, pero con sesgo bajista. \\
\texttt{bull\_call\_debit\_spread} & Dos calls, long strike menor, short strike
mayor, misma expiraci\'on, net debit y max loss definido. \\
\texttt{bear\_put\_debit\_spread} & Dos puts, long strike mayor, short strike
menor, misma expiraci\'on, net debit y max loss definido. \\
\texttt{protective\_put} & Requiere posici\'on subyacente o marca como hedge;
advierte si el costo es excesivo. \\
\texttt{collar} & Requiere subyacente, long put, short call, strikes coherentes
con spot y costo/cr\'edito calculable. \\
\bottomrule
\end{longtable}

Cada check puede devolver:

\begin{lstlisting}
{
  hardStops: string[],
  caps: Array<"watchlist" | "approved_small">,
  warnings: string[],
  scoreAdjustments: number
}
\end{lstlisting}

% ==========================================================
\section{Yahoo Finance Sentiment}
% ==========================================================

La capa de sentimiento de Reddit queda fuera del flujo principal y se conserva
solo como c\'odigo legado aislado. La nueva capa opcional es
\textbf{Yahoo Finance Sentiment}.

\subsection{Caracter\'isticas}

\begin{itemize}[noitemsep]
  \item Est\'a desactivada por defecto.
  \item No usa LLM externo ni API de pago.
  \item Usa endpoints p\'ublicos no oficiales de Yahoo Finance en modo
        best-effort.
  \item Si falla la recolecci\'on, devuelve baja se\~nal y no rompe la evaluaci\'on.
  \item Guarda documentos normalizados en SQLite mediante
        \texttt{YahooFinanceDocument}.
\end{itemize}

\subsection{Clasificaci\'on}

El clasificador local calcula:

\begin{itemize}[noitemsep]
  \item relevancia frente al ticker y catalizador;
  \item sentimiento bullish/bearish;
  \item hype;
  \item miedo;
  \item menciones de catalizador;
  \item preocupaciones legales/regulatorias.
\end{itemize}

Las etiquetas posibles son:

\[
\{\texttt{strong\_bullish},\texttt{bullish},\texttt{mixed},
\texttt{bearish},\texttt{strong\_bearish},\texttt{low\_signal}\}.
\]

\subsection{Overlay}

Si la confianza es alta y el sentimiento coincide con la direcci\'on del trade,
puede sumar $+3$ puntos. Si contradice la direcci\'on, puede restar $-3$.

Noticias severamente negativas, miedo alto o menciones legales/regulatorias
pueden limitar la decisi\'on a \texttt{watchlist} y reducir el multiplicador de
riesgo a 0.75.

% ==========================================================
\section{Base de datos y settings}
% ==========================================================

\subsection{Nuevos campos en \texttt{RiskSettings}}

\begin{lstlisting}
enableYahooFinanceSentiment Boolean @default(false)
sentimentLookbackDays       Int     @default(14)
sentimentMaxDocuments       Int     @default(10)
eventRiskEnabled            Boolean @default(true)
\end{lstlisting}

\subsection{Nuevo modelo \texttt{YahooFinanceDocument}}

\begin{lstlisting}
model YahooFinanceDocument {
  id             String   @id @default(cuid())
  portfolioId    String?
  ticker         String
  title          String
  publisher      String?
  summary        String?
  url            String?
  publishedAt    DateTime?
  collectedAt    DateTime @default(now())
  sentimentLabel String?
  sentimentScore Float?
  relevanceScore Float?
  hypeScore      Float?
  fearScore      Float?
  catalysts      String?
  concerns       String?

  @@index([ticker])
  @@index([publishedAt])
  @@index([collectedAt])
}
\end{lstlisting}

% ==========================================================
\section{Cambios visibles en UI}
% ==========================================================

\subsection{Trade Ideas}

La vista de detalle de ideas ahora muestra:

\begin{itemize}[noitemsep]
  \item PoP, fuente, confianza y advertencias.
  \item Kelly activado/desactivado y tope por Kelly fraccional.
  \item riesgo por eventos: nivel, earnings, macro, eventos y multiplicador;
  \item checks de estrategia: hard stops, caps, warnings y ajuste de score;
  \item Yahoo Finance Sentiment: etiqueta, confianza, relevancia, resumen,
        documentos usados, menciones y multiplicador.
\end{itemize}

\subsection{Risk Model}

La p\'agina de modelo ahora incluye una tabla de r\'egimen compuesto con SPY, QQQ,
breadth y VIX, mostrando peso, score, estado, raz\'on, confianza y advertencias.

\subsection{Settings}

La configuraci\'on muestra:

\begin{itemize}[noitemsep]
  \item activar/desactivar Yahoo Finance Sentiment;
  \item d\'ias de lookback;
  \item m\'aximo de documentos;
  \item activar/desactivar overlay de event risk.
\end{itemize}

% ==========================================================
\section{F\'ormula final de dimensionamiento}
% ==========================================================

El riesgo por multiplicadores es:

\[
R_{\text{mult}} =
B \cdot C_{\text{score}} \cdot L \cdot G \cdot V \cdot H \cdot M \cdot S \cdot E,
\]

donde:

\begin{longtable}{p{0.20\linewidth} p{0.68\linewidth}}
\toprule
\textbf{S\'imbolo} & \textbf{Descripci\'on} \\
\midrule
$B$ & Riesgo base por trade. \\
$C_{\text{score}}$ & Multiplicador por confianza del score. \\
$L$ & Liquidez de opciones. \\
$G$ & R\'egimen compuesto alineado con direcci\'on. \\
$V$ & Volatilidad/IV rank. \\
$H$ & Concentraci\'on sectorial o tem\'atica. \\
$M$ & Monte Carlo. \\
$S$ & Yahoo Finance Sentiment. \\
$E$ & Riesgo por eventos. \\
\bottomrule
\end{longtable}

Luego:

\[
R_{\text{final}} =
\min(R_{\text{mult}}, R_{\text{Kelly}}, R_{\text{sleeve}}, R_{\text{monthly}}, R_{\text{open}}).
\]

Si la decisi\'on final no empieza con \texttt{approved}, el riesgo sugerido es
cero.

% ==========================================================
\section{Comandos locales}
% ==========================================================

Para una base existente:

\begin{lstlisting}
npm install
npm run prisma:generate
npm run db:migrate:yahoo-overlays
npm run typecheck
npm test
npm run lint
npm run build
\end{lstlisting}

Para una base nueva:

\begin{lstlisting}
npm install
npm run db:init
npm run prisma:generate
npm run dev
\end{lstlisting}

No se agregaron variables de entorno nuevas. Alpaca sigue siendo opcional y de
solo lectura para snapshots de contratos de opciones.

% ==========================================================
\section{Limitaciones conocidas}
% ==========================================================

\begin{itemize}[noitemsep]
  \item Yahoo Finance es una fuente p\'ublica no oficial; puede fallar o cambiar.
  \item Earnings solo se eval\'ua cuando se proveen eventos o se agregue una fuente
        local/confiable.
  \item El calendario macro est\'atico debe mantenerse manualmente actualizado.
  \item La PoP por delta o IV proxy es menos confiable que Monte Carlo o IV real.
  \item El sistema no predice precios, no da asesor\'ia financiera y no coloca
        \'ordenes.
\end{itemize}

% ==========================================================
\section{Criterios de validaci\'on implementados}
% ==========================================================

Se agregaron pruebas unitarias para:

\begin{itemize}[noitemsep]
  \item r\'egimen compuesto bullish y datos faltantes;
  \item jerarqu\'ia PoP: Monte Carlo, delta proxy y no disponible;
  \item earnings crossing y hard stop por naked short;
  \item bull call debit spread v\'alido e inv\'alido;
  \item overlay de Yahoo Finance con preocupaci\'on legal/regulatoria.
\end{itemize}

La validaci\'on local pas\'o con:

\begin{lstlisting}
npm run prisma:generate
npm run typecheck
npm test
npm run lint
npm run build
\end{lstlisting}

\end{document}
